\subsection{Zakres pracy}

\paragraph{}Zakres pracy obejmuje analizę problemu rozpoznawania twarzy na podstawie obrazów 2D w warunkach częściowej okluzji, ze szczególnym uwzględnieniem syntetycznego zasłonięcia obszaru oczu. Badania przeprowadzono z wykorzystaniem wybranego zbioru danych zawierającego obrazy twarzy poddane procedurze wyrównania (alignment) oraz kontrolowanej augmentacji w postaci okluzji.

\paragraph{}W ramach projektu wykonano:

* porównanie wybranych modeli typu state-of-the-art (ArcFace Large, ArcFace Small, VGGFace) z autorskimi implementacjami opartymi na architekturze IResNet50,

* analizę wpływu rozproszonej uwagi CBAM (integracja w blokach rezydualnych) na skuteczność identyfikacji,

* analizę wpływu pomocniczej gałęzi Transformer na odporność na okluzję,

* ocenę skuteczności strategii transfer learningu i fine-tuningu w kontekście odporności na okluzję,

* ewaluację modeli w zadaniu identyfikacji tożsamości z wykorzystaniem metryk Rank-1 oraz Rank-3.

Zakres pracy ograniczono do porównania modelu bazowego oraz jego modyfikowanych wariantów architektonicznych. Nie przeprowadzono analizy alternatywnych rodzin modeli (np. w pełni transformatorowych), metod wymagających dodatkowych modalności biometrycznych ani zaawansowanych technik rekonstrukcji okluzji.

\paragraph{}Projekt nie obejmuje analizy danych wideo, danych trójwymiarowych (3D), zagadnień związanych z detekcją twarzy, ani optymalizacji modeli pod kątem wdrożeń produkcyjnych w systemach czasu rzeczywistego. Detekcja twarzy została wykorzystana wyłącznie jako etap wstępnego wyrównania danych wejściowych i nie stanowi przedmiotu analizy ani porównań w niniejszej pracy. Ponadto badania ograniczono do eksperymentów z wykorzystaniem syntetycznie generowanej okluzji, bez walidacji na rzeczywistych zbiorach danych zawierających naturalne przesłonięcia twarzy.